% Source: physical PDF pp. 333--336 (printed pp. 310--313), from the
% Section 30.2 heading through the final sentence immediately before
% Section 30.3 on physical p. 336.
% Inventory: Eqs. (30.2.1)--(30.2.20), five unnumbered display groups,
% no citations, figures, tables, footnotes, or typographic dividers.
% Corrected a summed index, a missing conjugate integration measure, and
% mismatched Lorentz indices in the electromagnetic propagator example.

\subsection{Superpropagators}

Usually the propagator can be obtained directly from the part of the action
of second order in the fields. If we write this quadratic part for a complex
field \(\phi_i\) (where \(i\) is a compound index including spacetime
coordinates as well as spin and species labels) in the form
\begin{equation}
  I_{\mathrm{quad}}[\phi]
  =-\sum_{ij}D_{ij}\phi_i^*\phi_j\,,
  \tag{30.2.1}\label{eq:30.2.1}
\end{equation}
with \(D_{ij}\) Hermitian, then as explained in Section 9.4 the propagator is
simply \(\Delta=D^{-1}\). The problem with this arises when the quadratic
part of the action is invariant under a (linearized) gauge transformation
\begin{equation}
  \phi_i\longrightarrow\phi_i+\xi_i\,,
  \tag{30.2.2}\label{eq:30.2.2}
\end{equation}
for some class of `vectors' \(\xi\). In this case we have
\begin{equation}
  \sum_jD_{ij}\xi_j=0\,,
  \tag{30.2.3}\label{eq:30.2.3}
\end{equation}
and we clearly cannot invert the `matrix' \(D_{ij}\). In electrodynamics the
problem arises because the Lagrangian density is invariant under a gauge
transformation \(A_\mu\longrightarrow A_\mu+\partial_\mu\Lambda\). We have
this problem here too: since the action is actually a functional of
\(\bar D^2S_n\) and not of \(S_n\) itself, it is invariant under the
transformation
\begin{equation}
  S_n\longrightarrow S_n+\bar D_{\dot\alpha}X_n^{\dot\alpha}\,,
  \tag{30.2.4}\label{eq:30.2.4}
\end{equation}
for any superfields \(X_n\).

In the electrodynamics of charged particles the problems introduced by gauge
invariance are typically handled by choosing a gauge, for instance by the
Faddeev--Popov--de Witt method described in Section 15.5. But the problems
raised here by invariance under the transformation \eqref{eq:30.2.2} are much
more like the problems of the effective field theory of photons at energies
below the threshold for producing charged particles, where the theory is
gauge-invariant simply because the action only involves gauge-invariant
fields. In such theories there is a simpler option. In addition to the
eigenvector \(\xi_i\) of \(D_{ij}\) with eigenvalue zero (taken for simplicity
to be in a unique direction) we can find a set of orthonormal eigenvectors
\(u_{\nu i}\) with eigenvalues \(d_\nu\neq0\):
\begin{equation}
  \sum_jD_{ij}u_{\nu j}=d_\nu u_{\nu i}\,,
  \qquad
  \sum_i u_{\nu i}^*u_{\nu'i}=\delta_{\nu\nu'}\,,
  \qquad
  \sum_i u_{\nu i}^*\xi_i=0\,.
  \tag{30.2.5}\label{eq:30.2.5}
\end{equation}
We can introduce a new set of integration variables \(\phi'\) and
\(\phi'_\nu\):
\begin{equation}
  \phi_i=\phi'\xi_i+\sum_\nu\phi'_\nu u_{\nu i}\,.
  \tag{30.2.6}\label{eq:30.2.6}
\end{equation}
The sort of integral encountered in Feynman diagram calculations of quantum
expectation values may then be written
\begin{align}
  &\int\left[\prod_i d\phi_i\,d\phi_i^*\right]
    \exp\{iI_{\mathrm{quad}}[\phi]\}
    \phi_a\cdots\phi_b^*\cdots
  \notag\\
  &\quad={}
    \mathcal{J}\int d\phi'\,d\phi'{}^*
    \int\left[\prod_\nu d\phi'_\nu\,d\phi_\nu'{}^*\right]
    \exp\left\{-i\sum_\nu d_\nu|\phi'_\nu|^2\right\}
  \notag\\
  &\qquad\cdot
    \left[\phi'\xi_a+\sum_\nu\phi'_\nu u_{\nu a}\right]\cdots
    \left[\phi'\xi_b+\sum_\nu\phi'_\nu u_{\nu b}\right]^*\cdots\,,
  \tag{30.2.7}\label{eq:30.2.7}
\end{align}
where \(\mathcal{J}\) is the Jacobian of the transformation
\eqref{eq:30.2.6}. The integral over \(\phi'\) and \(\phi'{}^*\) is of course
ill defined, because they do not appear in the argument of the exponential.
But this doesn't matter if the action only involves gauge-invariant fields,
because then \(\phi_a,\phi_b\), etc. will be contracted with `currents'
\(J_a,J_b\), etc. for which
\begin{equation}
  \sum_a\xi_aJ_a=0\,.
  \tag{30.2.8}\label{eq:30.2.8}
\end{equation}
We may thus write Eq.~\eqref{eq:30.2.7} as
\begin{align}
  &\int\left[\prod_i d\phi_i\,d\phi_i^*\right]
    \exp\{iI_{\mathrm{quad}}[\phi]\}
    \phi_a\cdots\phi_b^*\cdots
  \notag\\
  &\quad={}
    \mathcal{C}\int
    \left[\prod_\nu d\phi'_\nu\,d\phi_\nu'{}^*\right]
    \exp\left\{-i\sum_\nu d_\nu|\phi'_\nu|^2\right\}
    \left[\sum_\nu\phi'_\nu u_{\nu a}\right]\cdots
    \left[\sum_\nu\phi'_\nu u_{\nu b}\right]^*
    \cdots
  \notag\\
  &\qquad+\xi\text{-terms}\,,
  \tag{30.2.9}\label{eq:30.2.9}
\end{align}
where `\(\xi\)-terms' denotes terms proportional to one or more factors of
\(\xi_a,\xi_b\), etc., which vanish when contracted with \(J\)s satisfying
Eq.~\eqref{eq:30.2.8}, and \(\mathcal{C}\) is the infinite constant
\(\mathcal{J}\int d\phi'\,d\phi'{}^*\). The integral over the \(\phi'_\nu\) then gives
\begin{align}
  &\int\left[\prod_i d\phi_i\,d\phi_i^*\right]
    \exp\{iI_{\mathrm{quad}}[\phi]\}
    \phi_a\cdots\phi_b^*\cdots
  \notag\\
  &\qquad\propto
    \sum_{\text{pairings}}\bigl[-i\Delta_{ab}\bigr]\cdots
    +\xi\text{-terms}\,,
  \tag{30.2.10}\label{eq:30.2.10}
\end{align}
where the sum over pairings is over all ways of pairing indices on the
\(\phi\)s with indices on the \(\phi^*\)s, and \(\Delta_{ab}\) is the
propagator
\begin{equation}
  \Delta_{ab}=\sum_\nu\frac{u_{\nu a}u_{\nu b}^*}{d_\nu}\,.
  \tag{30.2.11}\label{eq:30.2.11}
\end{equation}
Instead of actually evaluating the sum \eqref{eq:30.2.11}, we can use its
defining property that
\begin{equation}
  \sum_cD_{ac}\Delta_{cb}
  =\sum_\nu u_{\nu a}u_{\nu b}^*
  \equiv\Pi_{ab}\,,
  \tag{30.2.12}\label{eq:30.2.12}
\end{equation}
with \(\Pi\) the projection operator on the space orthogonal to \(\xi\):
\begin{equation}
  \Pi^2=\Pi\,,
  \qquad
  \Pi\xi=0\,.
  \tag{30.2.13}\label{eq:30.2.13}
\end{equation}
The solution of Eq.~\eqref{eq:30.2.12} is unique only up to \(\xi\)-terms,
but these do not matter if the fields \(\phi_i\) appear in the action only in
gauge-invariant combinations.

For instance, in electrodynamics, we can write the kinematic part of the
action as
\[
  I_{\mathrm{quad}}[A]
  =-\frac14\int d^4x\,f_{\mu\nu}f^{\mu\nu}
  =+\frac12\int d^4x\,
    A^\mu\bigl(\Box\delta_\mu^\nu-\partial_\mu\partial^\nu\bigr)A_\nu\,.
\]
The differential operator
\(-\Box\delta_\mu^\nu+\partial^\nu\partial_\mu\) is not invertible, because
it has a zero eigenvalue with eigenvectors of the form
\(\xi^\mu=\partial^\mu\Lambda\). The projection matrix onto the space
orthogonal to these vectors is
\[
  \Pi_\mu{}^\nu(x,y)
  =\bigl[\delta_\mu^\nu-\partial_\mu\partial^\nu\Box^{-1}\bigr]
    \delta^4(x-y)\,,
\]
where \(\Box^{-1}\delta^4(x-y)\) is any solution of the equation
\(\Box[\Box^{-1}\delta^4(x-y)]=\delta^4(x-y)\). The defining equation for
the propagator is
\[
  \bigl(-\Box\delta_\mu^\lambda+\partial_\mu\partial^\lambda\bigr)
  \Delta_\lambda{}^\nu(x,y)
  =\Pi_\mu{}^\nu(x,y)\,,
\]
with the solution
\[
  \Delta_\lambda{}^\nu(x,y)
  =\delta_\lambda^\nu\Delta_F(x-y)
   +\partial_\lambda\partial^\nu\text{-terms}\,,
\]
where \(\Delta_F(x-y)\) is the usual Feynman propagator
\eqref{eq:6.2.16}, satisfying
\(\Box\Delta_F(x-y)=-\delta^4(x-y)\). (The factor \(1/2\) in the action
does not appear in the defining equation for the propagator here because
\(A_\mu\) is a real field. The origin of the \(-i\epsilon\) in the
denominator of the Fourier integral in Eq.~\eqref{eq:6.2.16} is explained
in the path-integral formalism in Section 9.2.)

Inspection of the first term in Eq.~\eqref{eq:30.1.9} shows that the
defining equation for the superpropagator of the potential superfields is
\begin{equation}
  \frac14D^2\bar D^2
  \Delta^S_{nm}(x,\theta,\bar\theta;x',\theta',\bar\theta')
  =\mathcal{P}\delta^4(x-x')\delta^2(\theta-\theta')
    \delta^2(\bar\theta-\bar\theta')\delta_{nm}\,,
  \tag{30.2.14}\label{eq:30.2.14}
\end{equation}
where \(\mathcal{P}\) is a superspace differential operator satisfying the
conditions for a projection operator
\begin{equation}
  \mathcal{P}^2=\mathcal{P}\,,
  \qquad
  \mathcal{P}\bar D_{\dot\alpha}=0\,,
  \tag{30.2.15}\label{eq:30.2.15}
\end{equation}
and
\(\delta^2(\theta-\theta')\delta^2(\bar\theta-\bar\theta')\)
is the fermionic delta function introduced in Eq.~\eqref{eq:26.6.8}.
The solution is
\begin{equation}
  \mathcal{P}=-\frac{1}{16\Box}D^2\bar D^2\,.
  \tag{30.2.16}\label{eq:30.2.16}
\end{equation}
(It is obvious that \(\mathcal{P}\bar D_{\dot\alpha}=0\). To check that
\(\mathcal{P}^2=\mathcal{P}\), we need to use Eq.~\eqref{eq:26.6.12}, which
shows that
\(\bar D^2D^2\bar D^2=-16\Box\bar D^2\).) The solution of
Eq.~\eqref{eq:30.2.14} is
\begin{align}
  \Delta^S_{nm}(x,\theta,\bar\theta;x',\theta',\bar\theta')
  &=-\frac{1}{4\Box}
    \delta^4(x-x')\delta^2(\theta-\theta')
    \delta^2(\bar\theta-\bar\theta')\delta_{nm}
  \notag\\
  &=\frac14\Delta_F(x-x')\delta^2(\theta-\theta')
    \delta^2(\bar\theta-\bar\theta')\delta_{nm}
    +\bar D_{\dot\alpha}\text{-terms}\,.
  \tag{30.2.17}\label{eq:30.2.17}
\end{align}
This is the superpropagator for a line created by a potential superfield
\(S_m^*(x',\theta',\bar\theta')\) and destroyed by a potential
superfield \(S_n(x,\theta,\bar\theta)\), which we must use in evaluating
supergraphs for the action
\eqref{eq:30.1.9}. To make contact with ordinary propagators, it is also of
some interest to consider the superpropagator of the line created by a
left-chiral superfield
\(\Phi_m^*(x',\theta',\bar\theta')\) and destroyed by a
left-chiral superfield \(\Phi_n(x,\theta,\bar\theta)\). These chiral
superfields are obtained by operating on
\(S_n(x,\theta,\bar\theta)\) with \(\bar D^2\) and on
\(S_m^*(x',\theta',\bar\theta')\) with
\((\bar D^2)^\dagger=D^{\prime\,2}\), so the propagator for the
left-chiral superfields is
\begin{equation}
  \Delta^\Phi_{nm}(x,\theta,\bar\theta;x',\theta',\bar\theta')
  =\frac14\bar D^2D^{\prime\,2}
    \Delta_F(x-x')\delta^2(\theta-\theta')
    \delta^2(\bar\theta-\bar\theta')\delta_{nm}\,.
  \tag{30.2.18}\label{eq:30.2.18}
\end{equation}
For instance, the term in the superpropagator of zeroth order in \(\theta\)
and \(\theta'\) is
\begin{equation}
  \left.
  \left[\Delta^\Phi_{nm}
    (x,\theta,\bar\theta;x',\theta',\bar\theta')\right]
  \right|_{\theta=\bar\theta=\theta'=\bar\theta'=0}
  =\frac14
    \left(\frac{\partial}{\partial\bar\theta}\right)^2
    \left(\frac{\partial}{\partial\theta'}\right)^2
    \Delta_F(x-x')\delta^2(\theta-\theta')
    \delta^2(\bar\theta-\bar\theta')\delta_{nm}\,.
  \tag{30.2.19}\label{eq:30.2.19}
\end{equation}
To evaluate this, we recall Eq.~\eqref{eq:26.6.8} for the fermionic delta
function:
\[
  \delta^2(\theta-\theta')\delta^2(\bar\theta-\bar\theta')
  =\frac14
    (\theta-\theta')^2(\bar\theta-\bar\theta')^2\,,
\]
from which we find
\((\partial/\partial\bar\theta)^2
(\partial/\partial\theta')^2
\delta^2(\theta-\theta')\delta^2(\bar\theta-\bar\theta')=4\).
Eq.~\eqref{eq:30.2.19} therefore yields
\begin{equation}
  \left.
  \left[\Delta^\Phi_{nm}
    (x,\theta,\bar\theta;x',\theta',\bar\theta')\right]
  \right|_{\theta=\bar\theta=\theta'=\bar\theta'=0}
  =\Delta_F(x-x')\delta_{nm}\,,
  \tag{30.2.20}\label{eq:30.2.20}
\end{equation}
which is just the usual propagator for the scalar components of the
superfield.
